
\section{Extension to general orthogonal alignments}
\label{app:general_W}

The main text (Theorem~\ref{thm:unified_bound}) states the bound at the identity alignment $W = I$. Here we lift Theorem~\ref{thm:unified_bound} to arbitrary orthogonal $W$, characterize the $W = I$ trace and $W = W_N$ nuclear specializations (the latter used in the Sec.~\ref{subsec:ablation} ablation), and treat the joint optimization over feature and head terms. The two specializations encode different priorities: $W_N$ minimizes the feature alignment residual but generally inflates the head term, while $W = I$ keeps the head term at its minimum (vanishing exactly under frozen \texttt{lm\_head}) and is SVD-free differentiable, properties that the LLM settings studied here---PTQ with FP16 head and frozen-head LoRA---require.

\subsection{The general bound for arbitrary $W$}
\label{app:general_bound}

The exact decomposition derived in Appendix~\ref{app:procrustes} (Eq.~\ref{eq:exact_decomposition_general}) holds for every $W \in \mathcal{O}(d)$ via the alignment-dependent similarity $\Omega_W$ of Eq.~\ref{eq:omega_def}, yielding a family of risk bounds parameterized by $W$:
\begin{equation}
|\mathcal{R}_T - \mathcal{R}_P| \;\le\; \underbrace{K_{\mathrm{feat}} \sqrt{(\rho_T - \rho_P)^2 + 2\rho_T\rho_P(1 - \Omega_W)}}_{\delta(W)} + \underbrace{K_{\mathrm{pred}} \|\Sigma_P^{1/2}(W H_T - H_P)\|_F}_{\gamma(W)}.
\end{equation}
The Procrustes solution $W_N$ minimizes $\delta(W)$ alone, yielding $\Omega_{W_N} = \Omega_N$ (the nuclear-norm similarity, Eq.~\ref{eq:omega_nuclear}). However, $W_N$ enters the head term $\gamma(W_N)$, which can be inflated when the rotation is applied to a nearly preserved head ($H_T \approx H_P$). This creates a fundamental trade-off: $W_N$ gives the tightest feature bound but may worsen the head bound, while $W = I$ preserves head coherence at the cost of a slightly looser feature bound.

\subsection{The trace specialization ($W = I$)}
\label{app:trace_specialization}

Setting $W = I$ recovers the main-text trace similarity $\Omega = \Omega_{W=I}$ (Eq.~\ref{eq:omega_def}). Two structural properties characterize this specialization.

\paragraph{Ordering of similarities.}
Since the Procrustes solution maximizes $\operatorname{Tr}(Z_T^\top Z_P W)$ over $\mathcal{O}(d)$:
\begin{equation}
\operatorname{Tr}(Z_T^\top Z_P) \;\le\; \|Z_T^\top Z_P\|_*,
\end{equation}
which gives $\Omega \le \Omega_N$. Equality holds if and only if $W_N = I$, i.e., $Z_T^\top Z_P$ is symmetric positive semidefinite.

\paragraph{When does $\Omega = \Omega_N$?}
Let $Z_T^\top Z_P = U\Sigma V^\top$ be the SVD. Then $\Omega = \Omega_N$ iff $VU^\top = I$, i.e., $U = V$, which holds when $Z_T^\top Z_P$ is symmetric positive semidefinite (SPSD). In post-training quantization with mild noise, $Z_T^\top Z_P \approx Z_T^\top Z_T$ is approximately SPSD, so $\Omega \approx \Omega_N$. As quantization becomes more aggressive, asymmetric perturbations cause $\Omega$ to fall below $\Omega_N$, and this gap reflects genuine asymmetric distortion that a rotation-invariant metric would mask.

\paragraph{Head-side simplification.}
At $W = I$, the head term $K_{\mathrm{pred}}\,\|\Sigma_P^{1/2}(H_T - H_P)\|_F$ carries no rotational contamination and vanishes whenever the prediction head is preserved (FP16 head under PTQ, frozen head under LoRA), recovering the regimes used throughout Sec.~\ref{sec:experiments} and the regularizer of Sec.~\ref{subsec:shape_reg}.


\subsection{Joint optimization on the Stiefel manifold}
\label{app:joint_opt}

The tightest bound over all orthogonal alignments requires solving:
\begin{equation}
\label{eq:joint_obj}
W_{\mathrm{opt}} = \arg\min_{W \in \mathcal{O}(d)}
\Big[
\underbrace{K_{\mathrm{feat}} \sqrt{\tfrac{1}{n}\|Z_T - Z_P W\|_F^2}}_{\delta(W)}
\;+\;
\underbrace{K_{\mathrm{pred}}\,\|\Sigma_P^{1/2}(W H_T - H_P)\|_F}_{\gamma(W)}
\Big].
\end{equation}
Unlike the standard Procrustes problem (which minimizes $\delta$ alone and admits the closed-form $W_N = VU^\top$), the addition of $\gamma(W)$ introduces a quadratic dependence on $W$ through the head term, related to the \emph{Weighted Orthogonal Procrustes Problem} (WOPP)~\cite{gower2004procrustes}.

\paragraph{Why no closed-form solution exists.}
The feature term $\delta(W)$ is concave in $\operatorname{Tr}(Z_T^\top Z_P W)$, while $\gamma(W)$ is convex in $W$. Their sum is neither convex nor concave on $\mathcal{O}(d)$, precluding a simple variational characterization.

\paragraph{Special case: isotropic features.}
If $\Sigma_P \approx \lambda I$, the head term simplifies to $K_{\mathrm{pred}}\sqrt{\lambda}\,\|W H_T - H_P\|_F$, which still depends on $W$ through the cross-term $\operatorname{Tr}(H_P^\top W H_T)$ when $\|W H_T - H_P\|_F^2$ is expanded. The joint objective then minimizes $K_{\mathrm{feat}}\sqrt{\|Z_T - Z_P W\|_F^2/n} + K_{\mathrm{pred}}\sqrt{\lambda}\|W H_T - H_P\|_F$, and each term individually admits a closed-form Procrustes minimizer; the sum does not. The isotropic assumption is, in any case, unrealistic for LLM representations, which exhibit well-documented anisotropy~\cite{ethayarajh2019contextual}.

\paragraph{Numerical solution.}
For the general case, $W_{\mathrm{opt}}$ can be approximated via Riemannian gradient descent on the Stiefel manifold, but this entangles scale, shape, and head components, destroying the clean physical interpretation.

\paragraph{Practical irrelevance for this paper's regimes.}
In the two primary settings:
\begin{itemize}
    \item \textbf{Quantization (GPTQ/BnB: $H_T = H_P$; GGUF: $H_T \approx H_P$):} The Procrustes-optimal $W_N$ minimizes the feature term $\delta(W)$ alone but inflates the head term $\gamma(W)$ when the head is approximately preserved (cf.~Appendix~\ref{app:general_bound} trade-off). We adopt $W = I$ in main analyses, where $\gamma$ vanishes for FP16-head retention.
    \item \textbf{LoRA ($H_T = H_P$):} The head term becomes $K_{\mathrm{pred}}\|\Sigma_P^{1/2}(W - I) H_T\|_F$, which vanishes at $W = I$ (and is generically positive at $W \neq I$). We adopt $W = I$ throughout the LoRA experiments and the shape regularizer of Sec.~\ref{subsec:shape_reg} as it yields the cleanest decomposition: under the frozen LoRA head ($H_T = H_P$) the choice $W = I$ gives $\Sigma_P^{1/2}(I H_T - H_P) = 0$, so $\gamma = 0$; the scale arm $(\rho_T - \rho_P)^2$ is $W$-invariant (observed but not actively regularized, since LoRA primarily perturbs shape rather than scale; Sec.~\ref{subsec:decompose}); and the shape arm $1 - \Omega$ becomes the single differentiable training-time target. The Procrustes-optimal $W_N \neq I$ would tighten the shape arm marginally ($\Omega_N \ge \Omega$) at the cost of inflating $\gamma$ and breaking the at-source shape contraction the regularizer relies on.
\end{itemize}
The Riemannian formulation becomes relevant only when comparing models with \emph{both} rotated features and divergent heads (e.g., full-parameter SFT). We leave this setting to future work.
